\section{Related Work}
\label{sec:related}
\textbf{Selecting items.} Anchor Points~\citep{vivek2024anchor} clusters items by the responses of many models and evaluates one representative per cluster; tinyBenchmarks~\citep{maiapolo2024tiny} clusters IRT item embeddings and adds IRT-based estimators (p-IRT, gp-IRT); metabench~\citep{kipnis2025metabench} keeps the most IRT-informative items and predicts scores with additive models; \citet{li2025active} learn a policy that selects items adaptively and imputes the rest; \citet{pacchiardi2024100} train an ``assessor'' on a handful of reference instances of the new model and report that random reference instances work about as well as the more elaborate selectors they tried, which agrees with part of our findings.
\textbf{Prediction under shift.} \citet{zhang2025benchmark} compare eleven benchmark-prediction methods on 19 benchmarks, observe that a random sample followed by regression is competitive, that methods degrade when the new model is more accurate than the source models, and propose an estimator inspired by augmented inverse-propensity weighting.
Our study is complementary to that work (of which we read only the abstract, and we did not use its code): we use the (public) score matrices they release, but design our own factorial comparison (6 selectors $\times$ several estimators), define the shift regimes through the true accuracy ranking of the models, quantify unbiasedness and interval coverage by repeated subsampling, and study a cross-fitted finite-population difference estimator with a cohort-adaptive shrinkage step. We do not claim priority for the idea of correcting a model-based estimate by sample residuals: it is classical in survey sampling (model-assisted/GREG estimation, \citealt{sarndal1992model}) and underlies prediction-powered inference~\citep{angelopoulos2023ppi} and double machine learning~\citep{chernozhukov2018dml}; the novelty is the exact finite-population form with out-of-fold residuals for item sampling, the empirical-Bayes shrinkage across a cohort of models sharing the sample, and the systematic empirical comparison.
\textbf{Evaluation data.} Our data are item-level correctness matrices derived from the Open LLM Leaderboard~\citep{fourrier2024openllm} (including MMLU-Pro,~\citealt{wang2024mmlupro}) and HELM~\citep{liang2023helm}, as packaged by \citet{zhang2025benchmark}.
